% Replace the current Evaluation chapter in the thesis with this block.
% The values below are generated from Evaluation/llm_evaluation_results.csv.

\chapter{Evaluation}

\section{Evaluation objectives}

The evaluation has two distinct objectives. The first concerns technical
usability: records should be valid, fields should be populated, and provenance
should be retained. The second concerns factual extraction quality: a populated
field must be correct and supported by the source text. These objectives must
not be confused, since a complete record can still contain an incorrect or an
overly inferred value.

\section{Automatic coverage results}

The rule-based consolidated database contains 1,153 records. Its review table
measures field completion and an internal heuristic confidence score. The
results in Table~\ref{tab:coverage-results} demonstrate high field population,
but neither measure is a human-judged accuracy score nor a measure of LLM
performance.

\begin{table}[H]
\centering
\begin{tabular}{lrrr}
\toprule
Source & Records & Mean completion rate & Mean heuristic confidence \\
\midrule
CNAV & 42 & 0.99 & 0.58 \\
Fiches actions SIAGE & 60 & 0.94 & 0.53 \\
Fondation de France & 43 & 0.97 & 0.58 \\
HCFEA & 23 & 0.97 & 0.54 \\
IReSP Autonomie & 131 & 0.98 & 0.63 \\
VADA & 854 & 1.00 & 0.60 \\
\midrule
Overall & 1,153 & 0.99 & 0.59 \\
\bottomrule
\end{tabular}
\caption{Automatic coverage of the rule-based consolidated database.}
\label{tab:coverage-results}
\end{table}

Completion covers title, territory, scale, target audience, stakeholder, date,
theme, source, and five descriptive fields. The confidence figure derives from
rule-based candidate scores. It is useful for prioritising manual review, not as
a calibrated probability that a field is correct.

\section{Human reference evaluation protocol}

The factual quality of LLM-assisted extraction was evaluated through a manually
annotated reference sample. Sixty documents were selected using a reproducible
random draw, stratified by source: ten documents were drawn from each of the
six corpus sources. This equal allocation was chosen to ensure that small
sources were represented in the evaluation; it does not reproduce the source
distribution of the complete corpus.

For each selected document, the LLM output was compared with its cleaned source
text. Seven fields were examined: \texttt{territoire}, \texttt{echelle},
\texttt{public\_vise}, \texttt{nature\_initiative},
\texttt{porteur\_initiative}, \texttt{date}, and \texttt{thematique}.
For the controlled categorical fields \texttt{public\_vise},
\texttt{echelle}, and \texttt{nature\_initiative}, the annotator recorded a
reference value and the evaluation reports micro precision, recall, micro F1,
macro F1, and exact-match rate. For the remaining fields, each model output was
judged as \textit{correct}, \textit{partiel}, \textit{incorrect},
\textit{absence correcte}, or \textit{non \'{e}valuable}. A prediction was
considered correct only when it was explicitly supported by the source text.

The labels \textit{absence correcte} and \textit{non \'{e}valuable} describe
two different situations. \textit{Absence correcte} was assigned when the model
left a field empty and the source text did not provide enough information to fill
it reliably; this is treated as a successful prediction. In contrast,
\textit{non \'{e}valuable} was assigned when the reviewer could not establish a
reliable reference value because the source text was too ambiguous, incomplete, or
degraded. Non-evaluable cases are reported separately and excluded from the
denominator of the accepted rate.

The annotation was performed by one reviewer. Consequently, the results
measure agreement with this documented reference annotation, but do not provide
an inter-annotator agreement estimate. A second independent annotation of a
subset would be a useful extension.

\section{Human reference evaluation results}

Table~\ref{tab:llm-categorical-results} reports the results for the three
controlled categorical fields. All scores range from 0 to 1, with 1 indicating
perfect agreement with the reference annotation. Precision measures whether the
labels produced by the model were appropriate; recall measures whether the labels
expected from the source text were retrieved; and exact match requires the full
predicted value to match the reference value for a document.

The target-audience field was evaluated against the French controlled vocabulary
used in the extraction prompt: \textit{aînés}, \textit{aînés et autres},
\textit{professionnels}, \textit{autre}, and \textit{public non connu}. The model
achieved perfect scores in this sample for target audience and initiative type.
Territorial scale produced one non-exact prediction. In this case, the reference
label was \textit{nationale}, whereas the model returned an empty value. The model
therefore introduced no incorrect scale label, which explains the
precision of 1.000, but it missed one expected label, producing a recall of 0.987.

\begin{table}[H]
\centering
\begin{tabular}{lrrrrr}
\toprule
Field & Precision & Recall & F1 micro & F1 macro & Exact match \\
\midrule
Target audience & 1.000 & 1.000 & 1.000 & 1.000 & 1.000 \\
Initiative type & 1.000 & 1.000 & 1.000 & 1.000 & 1.000 \\
Territorial scale & 1.000 & 0.987 & 0.993 & 0.993 & 0.983 \\
\bottomrule
\end{tabular}
\caption{Human reference evaluation of controlled categorical metadata
(n=60 documents per field).}
\label{tab:llm-categorical-results}
\end{table}

Table~\ref{tab:llm-judgment-results} presents the human judgments for the
remaining fields. ``Evaluable'' is the number of documents for which the reviewer
could establish a reference value; it is therefore equal to 60 minus the number of
non-evaluable cases. ``Accepted'' combines \textit{correct} and \textit{absence
correcte} judgments. The accepted rate is calculated only on evaluable documents,
so it measures the model's factual reliability where the source permits a clear
assessment. ``Partial'' identifies a value that is broadly supported but incomplete
or overly vague, whereas ``incorrect'' identifies an unsupported or wrong value.

\begin{table}[H]
\centering
\begin{tabular}{lrrrrrr}
\toprule
Field & Evaluable & Accepted & Partial & Incorrect & Non-evaluable & Accepted rate \\
\midrule
Territory & 54 & 53 & 1 & 0 & 6 & 0.981 \\
Stakeholder & 38 & 37 & 0 & 1 & 22 & 0.974 \\
Date & 56 & 44 & 7 & 5 & 4 & 0.786 \\
Theme & 60 & 60 & 0 & 0 & 0 & 1.000 \\
\bottomrule
\end{tabular}
\caption{Human judgments of LLM metadata extraction. ``Accepted'' includes
correct predictions and correctly empty fields.}
\label{tab:llm-judgment-results}
\end{table}

\subsection{Exploratory source-level analysis}

Because the reference sample contains ten documents from each source, the same
annotations can be inspected by source to identify possible sensitivity to
documentary format. This analysis is exploratory: ten documents per source are
not sufficient to estimate a stable source-specific performance level, and the
results should therefore be read as diagnostic evidence rather than as a ranking
of sources.

Target audience, initiative type, and theme were fully accepted in all six
source subsets. Territorial scale was also fully correct for five sources. The
only exception was Fondation de France, for which one expected scale label was
missed, resulting in an F1 score of 0.957. Table~\ref{tab:llm-source-results}
focuses on the fields for which source-dependent variation is most informative.
For territory, holder, and date, the notation $a/b$ gives the number of accepted
predictions over evaluable documents; ``--'' indicates that no document was
evaluable for that field.

\begin{table}[H]
\centering
\small
\begin{tabular}{lrrrr}
\toprule
Source & Scale F1 & Territory & Holder & Date \\
\midrule
CNAV & 1.000 & 8/9 (0.889) & 10/10 (1.000) & 4/8 (0.500) \\
Fiches actions SIAGE & 1.000 & 10/10 (1.000) & 8/9 (0.889) & 8/10 (0.800) \\
Fondation de France & 0.957 & 8/8 (1.000) & 7/7 (1.000) & 8/8 (1.000) \\
HCFEA & 1.000 & 9/9 (1.000) & 2/2 (1.000) & 9/10 (0.900) \\
IReSP Autonomie & 1.000 & 8/8 (1.000) & -- & 9/10 (0.900) \\
RFVAA & 1.000 & 10/10 (1.000) & 10/10 (1.000) & 6/10 (0.600) \\
\bottomrule
\end{tabular}
\caption{Exploratory source-level results from the stratified reference sample.
Territory, holder, and date are reported as accepted predictions over evaluable
documents, followed by the accepted rate.}
\label{tab:llm-source-results}
\end{table}

The results suggest that date extraction is particularly dependent on source
format. The accepted date rate ranges from 0.500 for CNAV to 1.000 for
Fondation de France, with RFVAA also showing a lower rate of 0.600. This is
consistent with the different roles of temporal expressions in the sources:
they may describe an initiative, an event, a publication, or an update to a
document. The IReSP holder field was non-evaluable in all ten sampled records,
which reflects the frequent absence of a clearly identifiable initiative holder
in research-project descriptions rather than a measured model failure.

\section{Results analysis and limitations}

The human evaluation confirms that, for this stratified sample, the LLM can
produce reliable structured metadata for target audience, initiative type,
territorial scale, theme, and territory. The results should nevertheless be
interpreted as sample-level evidence rather than as a guarantee of identical
performance on every document in the corpus. In particular, the equal source
allocation makes the evaluation appropriate for detecting source-specific
weaknesses, but it is not an estimate weighted by the distribution of the
complete corpus.

Date extraction is the clearest weakness. Among the 56 evaluable documents,
only 44 predictions were fully accepted; seven were partial and five were
incorrect. This result is consistent with the heterogeneous role of dates in
the sources: a year can refer to an initiative start, publication, event,
funding period, or another temporal mention. Future versions of the schema
should distinguish the temporal expression from its semantic role, for example
by separating publication date, initiative start date, and event date.

The stakeholder field has a high accepted rate among evaluable documents
(0.974), but 22 of the 60 documents were non-evaluable. This does not
necessarily indicate a model error; rather, many source texts do not clearly
distinguish the initiative holder from a partner, funder, or organisation merely
mentioned in the document. The interface should therefore preserve uncertainty
and provenance instead of presenting every extracted stakeholder as an
established fact.

These findings complement, rather than replace, automatic coverage measures.
The completed reference evaluation establishes factual evidence for the seven
reviewed metadata fields, while the remaining free-text descriptions still
require separate qualitative validation before being used as factual claims.
The modest sample size and the absence of a second annotator remain limitations.
Future work should extend the reference set, measure inter-annotator agreement,
and evaluate whether prompt or model changes improve the date and stakeholder
fields.
